\section{变质量火箭的动力学模型}
\subsection{理想推进关系与质量比下界}
设质量为$m$，相对排气速度大小为$c$，排出质量为$-\mathrm dm>0$。微元动量平衡为
\begin{equation}m\mathrm dv=-c\mathrm dm+F_{\rm ext}\mathrm dt.\end{equation}
忽略外力且$c$恒定时，
\begin{equation}\Delta v_{\rm ideal}=c\ln(m_0/m_f),\qquad I_T=\int_0^{t_f}T\dd t=c(m_0-m_f).\label{eq:tsiolkovsky}\end{equation}
速度增量涉及$\int T/m\dd t$，不能以总冲量除以初始质量代替。

目标半径$r_*=R_E+h_*$、圆轨道速度$v_c=\sqrt{\mu/r_*}$。对从$r_0$静止出发且$r\ge r_0$的确定性轨迹，令
\begin{equation}\mathcal E=V^2/2+\mu(1/r_0-1/r).\end{equation}
无随机外力且阻力耗能时，$\dot{\mathcal E}\le(T/m)V\le(T/m)\sqrt{2\mathcal E}$。积分并代入圆轨道末态得到
\begin{equation}\frac{m_0}{m_f}\ge\exp\left[\frac1c\sqrt{\frac{2\mu}{r_0}-\frac{\mu}{r_*}}\right].\end{equation}
算例中该界约为6.245，仅用于能力筛查；它不是充分可行条件，也不对任意高斯外力样本成立。

\subsection{二维极坐标运动方程}
地球取不旋转球对称引力源，火箭取固定惯性平面内的变质量质点。忽略升力、分级与姿态执行机构动态。定义
\begin{equation}h=r-R_E,\quad u=\dot r,\quad w=r\dot\theta,\quad V=\sqrt{u^2+w^2},\quad\beta=\operatorname{atan2}(w,u).\end{equation}
控制为推力相对速度方向的偏角$\delta$及流率$q=-\dot m\ge0$，故$T=cq$、$\alpha=\beta+\delta$。起飞前5\,s保持$\alpha=0$，零速度附近采用方向保护。状态$x=(r,\theta,u,w,m,I_T)^{\mathsf T}$满足
\begin{equation}
\begin{aligned}
\dot r&=u,&\dot\theta&=w/r,\\
\dot u&=\frac{w^2}{r}+\frac{cq\cos\alpha}{m}-\frac{\mu}{r^2}-\frac{kVu}{m}+a_{\varepsilon,r},\\
\dot w&=-\frac{uw}{r}+\frac{cq\sin\alpha}{m}-\frac{kVw}{m}+a_{\varepsilon,\theta},\\
\dot m&=-q,&\dot I_T&=cq .
\end{aligned}\label{eq:dynamics}
\end{equation}
其中$\rho(h)=\rho_0e^{-\max(h,0)/H_a}$、$k=\rho C_DA_f/2$。阻力矢量为$-kV(u\bm e_r+w\bm e_\theta)$，大小为$kV^2$，在静止大气假设下与惯性速度反向。$w^2/r$与$-uw/r$为坐标运动学项，不是附加外力。优化允许主动偏角，不将全程视为严格零偏角重力转弯。

\subsection{密度相关随机激励}
\begin{equation}a_{\varepsilon,j}(t,h)=\sigma_{0,j}\frac{\rho(h)}{\rho_{\rm ref}}\eta_j(t),\quad j\in\{r,\theta\}.\label{eq:noise}\end{equation}
在间隔1\,s的网格上生成独立标准高斯样本$\xi_{j,\ell}$，作方差归一化插值：
\begin{equation}\eta_j(t)=\frac{(1-\lambda)\xi_{j,\ell}+\lambda\xi_{j,\ell+1}}{\sqrt{(1-\lambda)^2+\lambda^2}},\quad\lambda=t-\lfloor t\rfloor.\end{equation}
该过程具有有限时间相关性和单位边缘方差，可作为固定样本路径上的ODE激励。它不是理想白噪声；归一化因子使插值后过程也不是严格分段线性。扰动标准差与密度成正比，方差与密度平方成正比。
